% ============================================================================
% Rapport Technique — Drone de surface à voile sur T-Beam ESP32
% Auteurs : Mohamed EL JILY (code original), Facundo Arito (refonte)
% Version : 2 (mai 2026)
% Ce fichier est conçu pour être mis à jour au fil du projet.
% Pour ajouter une phase : créer un nouveau \section{} dans le fichier
%   phases/phaseN.tex et l'inclure avec \input{phases/phaseN}.
% ============================================================================

\documentclass[12pt,a4paper]{article}

% ── Encodage (pdflatex) ──────────────────────────────────────────────────────
\usepackage[utf8]{inputenc}
\usepackage[T1]{fontenc}
\usepackage{lmodern}
\usepackage{textcomp}

% ── Langue ───────────────────────────────────────────────────────────────────
% texlive-lang-french non installé : traductions manuelles.
\renewcommand{\contentsname}{Table des mati\`eres}
\renewcommand{\abstractname}{R\'esum\'e}
\renewcommand{\appendixname}{Annexe}
\renewcommand{\refname}{R\'ef\'erences}
\renewcommand{\figurename}{Figure}
\renewcommand{\tablename}{Tableau}

% ── Mise en page ─────────────────────────────────────────────────────────────
\usepackage[top=2.5cm, bottom=2.5cm, left=2.8cm, right=2.8cm]{geometry}
\usepackage{parskip}
\usepackage{microtype}
\usepackage{setspace}
\onehalfspacing

% ── Couleurs ─────────────────────────────────────────────────────────────────
\usepackage[dvipsnames]{xcolor}
\definecolor{codebg}{HTML}{1E1E2E}
\definecolor{codefg}{HTML}{CDD6F4}
\definecolor{commentfg}{HTML}{6C7086}
\definecolor{keywordfg}{HTML}{CBA6F7}
\definecolor{stringfg}{HTML}{A6E3A1}
\definecolor{numberfg}{HTML}{FAB387}
\definecolor{warnbg}{HTML}{FFF8E1}
\definecolor{warnborder}{HTML}{F59E0B}
\definecolor{okbg}{HTML}{F0FDF4}
\definecolor{okborder}{HTML}{16A34A}
\definecolor{infobg}{HTML}{EFF6FF}
\definecolor{infoborder}{HTML}{2563EB}
\definecolor{tableheadbg}{HTML}{374151}
\definecolor{tableheadfg}{HTML}{FFFFFF}
\definecolor{tablerowalt}{HTML}{F9FAFB}

% ── Boîtes colorées ──────────────────────────────────────────────────────────
\usepackage{tcolorbox}
\tcbuselibrary{skins, breakable}

\newtcolorbox{warningbox}{
  colback=warnbg, colframe=warnborder,
  fonttitle=\bfseries, title={\color{warnborder}$\blacktriangle$\ Avertissement},
  breakable, boxrule=1pt, left=6pt, right=6pt
}
\newtcolorbox{okbox}{
  colback=okbg, colframe=okborder,
  fonttitle=\bfseries, title={\color{okborder}$\checkmark$\ Résolu},
  breakable, boxrule=1pt, left=6pt, right=6pt
}
\newtcolorbox{infobox}[1][]{
  colback=infobg, colframe=infoborder,
  fonttitle=\bfseries, title={\color{infoborder}$\blacksquare$\ #1},
  breakable, boxrule=1pt, left=6pt, right=6pt
}

% ── Listings (code source) ────────────────────────────────────────────────────
\usepackage{listings}
\usepackage{lstautogobble}

% Les caractères non-ASCII dans les listings sont remplacés par des équivalents
% ASCII directement dans la source (voir ci-dessous). Pas de literate ici.

\lstdefinestyle{cpp}{
  language=C++,
  backgroundcolor=\color{codebg},
  basicstyle=\ttfamily\footnotesize\color{codefg},
  keywordstyle=\color{keywordfg}\bfseries,
  stringstyle=\color{stringfg},
  commentstyle=\color{commentfg}\itshape,
  numberstyle=\tiny\color{numberfg},
  numbers=left, numbersep=8pt,
  breaklines=true, breakatwhitespace=true,
  frame=single, rulecolor=\color{codebg!40!black},
  tabsize=4, showstringspaces=false,
  autogobble=true,
  morekeywords={uint8_t,uint16_t,uint32_t,int8_t,int16_t,int32_t,bool,
                constexpr,noexcept,nullptr,override,IRAM_ATTR,portMUX_TYPE},
}

\lstdefinestyle{shell}{
  language=bash,
  backgroundcolor=\color{codebg},
  basicstyle=\ttfamily\footnotesize\color{codefg},
  breaklines=true, frame=single, rulecolor=\color{codebg!40!black},
  showstringspaces=false,
}

\lstdefinestyle{json}{
  backgroundcolor=\color{codebg},
  basicstyle=\ttfamily\footnotesize\color{codefg},
  breaklines=true, frame=single, rulecolor=\color{codebg!40!black},
  showstringspaces=false,
}

% ── Tableaux ─────────────────────────────────────────────────────────────────
\usepackage{booktabs}
\usepackage{array}
\usepackage{longtable}
\usepackage{colortbl}
\usepackage{multirow}

\newcolumntype{L}[1]{>{\raggedright\arraybackslash}p{#1}}
\newcolumntype{C}[1]{>{\centering\arraybackslash}p{#1}}

% ── Hyperliens ───────────────────────────────────────────────────────────────
\usepackage[hidelinks]{hyperref}
\hypersetup{
  pdftitle={Rapport Technique -- Drone de surface à voile},
  pdfauthor={Facundo Arito},
  bookmarksnumbered=true,
}

% ── Mise en forme des titres ──────────────────────────────────────────────────
\usepackage{titlesec}
\titleformat{\section}{\Large\bfseries\color{tableheadbg}}{\thesection.}{0.7em}{}[\vspace{-0.3em}\color{infoborder}\rule{\linewidth}{1.5pt}]
\titleformat{\subsection}{\large\bfseries\color{tableheadbg!80!black}}{\thesubsection.}{0.7em}{}
\titleformat{\subsubsection}{\normalsize\bfseries}{\thesubsubsection.}{0.6em}{}

% ── Listes ───────────────────────────────────────────────────────────────────
\usepackage{enumitem}
\setlist{itemsep=2pt, topsep=4pt}

% ── En-têtes et pieds de page ─────────────────────────────────────────────────
\usepackage{fancyhdr}
\setlength{\headheight}{15pt}
\pagestyle{fancy}
\fancyhf{}
\fancyhead[L]{\small Drone de surface à voile — Rapport technique}
\fancyhead[R]{\small v2 — \today}
\fancyfoot[C]{\small\thepage}
\renewcommand{\headrulewidth}{0.4pt}

% ── Utilitaires ──────────────────────────────────────────────────────────────
\usepackage{graphicx}
\usepackage{amsmath}
\usepackage{amssymb}

% ─────────────────────────────────────────────────────────────────────────────
%  Commandes utilitaires
% ─────────────────────────────────────────────────────────────────────────────
\newcommand{\gpio}[1]{\texttt{GPIO#1}}
\newcommand{\us}[1]{\texttt{#1\,µs}}
\newcommand{\file}[1]{\texttt{#1}}
\newcommand{\status}[1]{\textbf{[#1]}}
\newcommand{\done}{\textcolor{okborder}{\textbf{[OK]}}}
\newcommand{\todo}{\textcolor{warnborder}{\textbf{[ ]}}}

% ─────────────────────────────────────────────────────────────────────────────
%  DOCUMENT
% ─────────────────────────────────────────────────────────────────────────────
\begin{document}

% ── Page de titre ─────────────────────────────────────────────────────────────
\begin{titlepage}
  \centering
  \vspace*{2cm}
  {\Huge\bfseries Drone de surface à voile\par}
  \vspace{0.5cm}
  {\Large Rapport technique — Développement du firmware et de l'électronique\par}
  \vspace{2cm}
  \rule{\linewidth}{1pt}
  \vspace{0.5cm}

  \begin{tabular}{ll}
    \textbf{Plateforme :} & LilyGO T-Beam V1.1 (ESP32 + SX1276 + GPS u-blox + AXP192) \\[4pt]
    \textbf{Langage :}    & Arduino / C++, cible ESP32 Arduino Core 3.x \\[4pt]
    \textbf{Auteur (code original) :} & Mohamed EL JILY \\[4pt]
    \textbf{Auteur (refonte) :}       & Facundo Arito \\[4pt]
    \textbf{Version :}    & 2.0 — \today \\
  \end{tabular}

  \vspace{0.5cm}
  \rule{\linewidth}{1pt}
  \vspace{2cm}

  \begin{abstract}
    Ce document retrace l'intégralité du développement informatique et électronique
    du drone de surface à voile basé sur le module T-Beam d'Espressif/LilyGO.
    Il couvre le code original et ses limitations, la réécriture complète de
    l'architecture logicielle, les deux révisions du PCB, tous les problèmes
    matériels et logiciels rencontrés, ainsi que les solutions mises en place.
    Le document est structuré pour être mis à jour facilement au fil des phases
    futures du projet.
  \end{abstract}

  \vfill
  \textit{Projet de recherche — Drone autonome de surveillance océanographique}
\end{titlepage}

% ── Table des matières ────────────────────────────────────────────────────────
\tableofcontents
\newpage

% =============================================================================
\section{Présentation du projet}
% =============================================================================

\subsection{Concept physique}

Le drone est un véhicule de surface propulsé principalement par la voile. Sa
conception est originale : la voile est à profil d'aile et tourne librement
autour d'un axe vertical. L'aérodynamique passive oriente automatiquement le
profil à environ 45° de l'axe du drone par rapport au vent. Deux actionneurs
principaux permettent de naviguer~:

\begin{itemize}
  \item \textbf{Servo aileron (voile) — Futaba S3003~:} contrôle un petit volet
    sur le bord de fuite de la voile. Sa course utile est strictement de
    $\pm10^{\circ}$. La déflexion détermine dans quel sens la voile propulse le drone,
    mais n'est pas continue~: seules les deux positions extrêmes ont une
    signification physique (babord ou tribord).

  \item \textbf{Servo safran/rotor — Graupner Regatta ECO II (mod. 5176)~:}
    contrôle le différentiel entre la rotation libre de la voile et un rotor
    fixe à la poupe. C'est le principal actionneur de cap (lacet). C'est un
    \emph{treuil de voile multi-tour positonnel}~: le signal PWM commande une
    position cible de l'arbre (et non une vitesse). Centre = \us{1500},
    course complète = 6 tours, maintien de position sous charge.

  \item \textbf{Propulseurs (ESC)~:} utilisés uniquement en l'absence de vent
    suffisant. Permettent la propulsion et la direction par poussée différentielle.
\end{itemize}

\begin{infobox}[Principe de navigation]
  Sans capteur de vent, la direction du vent est \textbf{déduite} du track GPS
  et de l'état des actionneurs. L'estimation nécessite un déplacement minimal
  de 30~m et deux amures (bords) observées.
\end{infobox}

\subsection{Matériel principal}

\begin{longtable}{L{4cm} L{9cm}}
  \toprule
  \rowcolor{tableheadbg}
  \textcolor{tableheadfg}{\textbf{Composant}} &
  \textcolor{tableheadfg}{\textbf{Description}} \\
  \midrule
  \endhead
  \bottomrule
  \endfoot
  LilyGO T-Beam V1.1 & Module intégrant ESP32 WROVER, LoRa SX1276 433~MHz,
                        GPS u-blox NEO-6M, AXP192 (gestion d'énergie),
                        porte-batterie 18650 \\[3pt]
  Futaba S3003 & Servo aileron voile, centre à \us{1520}, $\pm10^{\circ}$ \\[3pt]
  Graupner Regatta ECO II & Treuil safran multi-tour, \us{1500} = centre,
                            1000--2000~µs = plage, jusqu'à 3,3~A en charge \\[3pt]
  Pro-Tronik PTR-6A & Émetteur 6 voies FHSS 2,4~GHz, PWM 1000--2000~µs \\[3pt]
  Pro-Tronik R8X & Récepteur 8 voies, PWM standard 50~Hz \\[3pt]
  Pro-Tronik Black Fet & ESC série, programmation par carte EPRG-3 \\[3pt]
  Antenne GPS & Taoglas ADFGP.25A, active, connecteur u.FL \\[3pt]
  PCB v1 / PCB v2 & Carte d'interface custom (KiCad 9) \\
\end{longtable}

% =============================================================================
\section{Code original — \texttt{boat.ino}}
% =============================================================================

\subsection{Architecture et fonctionnement}

Le code original (\file{Informatica/Arduino/boat/boat.ino}), écrit par Mohamed
EL JILY, constitue le point de départ du projet. Il est organisé en modules~:

\begin{itemize}
  \item \file{Gps.cpp/.hpp} — lecture GPS (Serial1, TinyGPSPlus), historique
    de positions, lissage du cap
  \item \file{ServoControl.cpp/.hpp} — contrôle servo via ESP32Servo
  \item \file{RadioReceiver.cpp/.hpp} — lecture PWM par interruptions GPIO
    (3 canaux~: PWM1, PWM2, SEL)
  \item \file{MotorControl.cpp/.hpp} — armement et contrôle ESC
  \item \file{boat.ino} — boucle principale, logique de navigation à la voile
    (observation du vent, waypoints, lofer/abattre/virement de bord)
  \item \file{config\_pins.h} — correspondance GPIO
\end{itemize}

La logique de navigation est implémentée~: calcul de l'angle relatif vent/cap,
changement d'amure automatique (VDB), lofer/abattre. Les waypoints sont reçus
par LoRa depuis un serveur en JSON.

\subsection{Problèmes et limitations identifiés}

Lors de la reprise du projet, plusieurs problèmes bloquants ont été identifiés~:

\subsubsection{Bibliothèque ESP32Servo — conflit LEDC/interruptions}

\begin{lstlisting}[style=cpp, caption={ServoControl.cpp original -- utilisation d'ESP32Servo}]
void ServoControl::init() {
    sailServo.setPeriodHertz(SERVO_PWM_FREQ);
    sailServo.attach(SERVO_SAIL_PIN, SERVO_PULSE_MIN, SERVO_PULSE_MAX);
    ...
}
void ServoControl::setSailAngle(float angle) {
    angle += 90;           // decalage vers 0-180deg
    sailServo.write(constrain((int)angle, 0, 180));
}
\end{lstlisting}

\begin{warningbox}
  ESP32Servo utilise le périphérique \textbf{LEDC} de l'ESP32.
  LEDC est incompatible avec le décodage PWM par interruptions GPIO~:
  lorsque LEDC est actif, les interruptions RC sont perturbées et les
  valeurs lues deviennent aléatoires. Ce conflit rend impossible le
  contrôle simultané des servos et la lecture du récepteur radio.
\end{warningbox}

\subsubsection{Attribution des broches RC — conflit I2C/AXP192}

\begin{lstlisting}[style=cpp, caption={config\_pins.h original -- broches RC sur I2C}]
#define RADIO_PWM1_IN 21   // GPIO 21 -- CH4 Aileron (= I2C SDA AXP192!)
#define RADIO_PWM2_IN 22   // GPIO 22 -- CH2 Safran  (= I2C SCL AXP192!)
#define RADIO_SEL_IN  23   // GPIO 23 -- CH6 Select. (= LoRa RST hardware!)
\end{lstlisting}

\begin{warningbox}
  GPIO21 et GPIO22 sont les broches I2C internes du T-Beam utilisées
  par l'AXP192 (gestionnaire d'énergie). Les utiliser comme entrées
  RC sans couper proprement l'I2C provoque des instabilités sur le bus.
  De plus, \gpio{23} est connecté en interne au RESET du SX1276 (LoRa)
  sur le T-Beam V1.1~: attacher une interruption sur ce pin peut
  provoquer des resets intempestifs du module radio.
\end{warningbox}

\subsubsection{Broche batterie GPIO35 — drain de 20~mA continu}

\begin{lstlisting}[style=cpp, caption={config\_pins.h original -- broche batterie}]
// GPIO35 (Battery ADC) - Voltage divider R~370 Ohm total
// Ne JAMAIS lire GPIO35 : drain 20 mA permanent depuis la batterie !
\end{lstlisting}

Le T-Beam V1.1 possède un diviseur résistif de \emph{faible impédance}
($\approx 370$~$\Omega$ total) soudé en permanence sur \gpio{35}. Même sans
lire l'ADC, ce diviseur consomme $\approx 20$~mA en continu, soit environ
1,8~W~! Le code original utilisait \gpio{15} à la place, ce qui est correct,
mais sans documenter clairement ce problème.

\subsubsection{Broche LoRa RST — erreur de documentation}

\begin{lstlisting}[style=cpp, caption={config\_pins.h original -- LoRa RST incorrect}]
#define LORA_RST 14        // GPIO 14 (INCORRECT -- hardware = GPIO 23)
\end{lstlisting}

Le hardware T-Beam V1.1 connecte le RESET du SX1276 à \gpio{23}, pas à
\gpio{14} comme documenté dans les exemples de LilyGO. Utiliser \gpio{14}
ne déclenche aucun reset matériel, mais peut causer des comportements
inattendus si ce pin est utilisé à d'autres fins.

\subsubsection{Gestion ESC — absence d'armement et de limitation de variation}

Le code original ne met pas en œuvre de procédure d'armement ESC~: le moteur
peut démarrer immédiatement si le signal de commande est trop élevé au
démarrage. De même, aucune limitation de taux de variation (\emph{slew rate})
n'est implémentée, ce qui peut provoquer des à-coups mécaniques lors de
changements de consigne brusques.

% =============================================================================
\section{PCB Version 1 — Carte d'interface originale}
% =============================================================================

\subsection{Description}

La première carte PCB (fichier \file{PCB\_Electronica.kicad\_pcb}) réalise
l'interface entre le T-Beam et les actionneurs. Elle comprend~:

\begin{itemize}
  \item Des connecteurs 3 broches pour les servos/ESC (PWM, +5V, GND)
  \item Des connecteurs pour le récepteur RC
  \item Un régulateur abaisseur 5~V (Pololu 2866) alimenté par la batterie
  \item Des bornes à vis pour la batterie et les sorties de puissance
\end{itemize}

\subsection{Problèmes identifiés sur le PCB V1}

\begin{longtable}{L{3.5cm} L{9.5cm}}
  \toprule
  \rowcolor{tableheadbg}
  \textcolor{tableheadfg}{\textbf{Problème}} &
  \textcolor{tableheadfg}{\textbf{Impact}} \\
  \midrule
  \endhead
  \bottomrule
  \endfoot
  Signaux PWM 3,3~V sans amplification niveau & Les servos et ESC
    (logique 5~V) peuvent ne pas détecter correctement les impulsions \\[3pt]
  Broches RC sur GPIO21/22 & Conflict I2C/AXP192, impossible d'utiliser
    l'I2C après l'init sans couper le bus \\[3pt]
  ESC sur GPIO32/33 & GPIO32/33 = LoRa DIO1/DIO2 sur le T-Beam~;
    non bloquant en phase 1--2 mais à surveiller \\[3pt]
  GPIO35 pour batterie & Diviseur 370~$\Omega$ = 20~mA continu, non utilisable \\[3pt]
  Pas de protection en sens inverse & Risque en cas d'alimentation incorrecte \\
\end{longtable}

% =============================================================================
\section{Réécriture logicielle — Architecture Post-Claude}
% =============================================================================

\subsection{Motivations}

La réécriture complète du firmware vise à corriger toutes les limitations
identifiées et à préparer l'architecture pour les phases autonomes futures.
Les objectifs sont~:

\begin{enumerate}
  \item Utiliser MCPWM au lieu de LEDC pour les sorties actionneurs
  \item Décodage RC entièrement par interruptions avec protection ISR
  \item Architecture OOP en couches, extensible par phases
  \item Armement ESC et limitation de variation (\emph{slew rate})
  \item Gestion propre de l'énergie via AXP192
  \item GPS, LoRa et batterie fonctionnant simultanément
\end{enumerate}

\subsection{Architecture en couches}

\begin{lstlisting}[style=shell, caption={Arborescence du projet Post-Claude}]
Informatica/Post-Claude/
  main/                    <- Sketch Arduino production
    main.ino
    src/
      app/DroneApp.cpp     <- Orchestrateur principal (scheduler)
      config/BoardConfig.h <- Affectation des GPIO
      config/Calibration.h <- Valeurs uss servo/ESC
      core/Types.h         <- Types partages (RcFrame, ActuatorCommand...)
      control/ModeManager  <- CH5 -> mode de vol
      control/ManualController <- Controle manuel RC
      drivers/RcReceiver   <- Lecture PWM par interruptions
      drivers/McpwmActuators <- Sorties MCPWM
      drivers/AxpPower     <- Init AXP192
      drivers/BatteryAdc   <- ADC tension batterie
      drivers/GpsUart      <- Lecture GPS (TinyGPSPlus)
      drivers/LoRaRadio    <- Driver SX1276
      comm/LoRaComm        <- Protocole JSON LoRa
      navigation/          <- Navigator, MissionPlan, MissionManager
  PCBv3/                   <- Sketch pour test sur platine d'essai
\end{lstlisting}

\subsection{Décisions techniques fondamentales}

\subsubsection{MCPWM obligatoire — jamais LEDC}

\begin{lstlisting}[style=cpp, caption={McpwmActuators.cpp -- initialisation MCPWM}]
bool McpwmActuators::begin() {
    mcpwm_gpio_init(MCPWM_UNIT_0, MCPWM0A, BoardConfig::SAIL_SERVO_PIN);
    mcpwm_gpio_init(MCPWM_UNIT_0, MCPWM0B, BoardConfig::ROTOR_SERVO_PIN);
    mcpwm_gpio_init(MCPWM_UNIT_0, MCPWM1A, BoardConfig::ESC1_PIN);
    mcpwm_config_t cfg = {};
    cfg.frequency    = Calibration::PWM_FREQ_HZ;   // 50 Hz
    cfg.duty_mode    = MCPWM_DUTY_MODE_0;           // actif haut
    cfg.counter_mode = MCPWM_UP_COUNTER;
    mcpwm_init(MCPWM_UNIT_0, MCPWM_TIMER_0, &cfg);
    mcpwm_init(MCPWM_UNIT_0, MCPWM_TIMER_1, &cfg);
    ...
}
\end{lstlisting}

Le périphérique MCPWM de l'ESP32 est indépendant des interruptions GPIO~:
les deux peuvent coexister sans perturbation mutuelle. Cette contrainte est
\textbf{non négociable} pour ce projet.

\subsubsection{Interruptions RC avec protection portMUX}

\begin{lstlisting}[style=cpp, caption={RcReceiver.cpp -- ISR et section critique}]
void IRAM_ATTR RcReceiver::handleEdge(uint8_t pin, Ch ch) {
    const uint32_t now = micros();
    const int      idx = static_cast<int>(ch);

    if (gpio_get_level(static_cast<gpio_num_t>(pin))) {
        portENTER_CRITICAL_ISR(&mux_);
        riseUs_[idx] = now;
        portEXIT_CRITICAL_ISR(&mux_);
    } else {
        uint32_t rise;
        portENTER_CRITICAL_ISR(&mux_);
        rise = riseUs_[idx];
        portEXIT_CRITICAL_ISR(&mux_);
        const uint32_t width = now - rise;
        if (width >= RC_VALID_MIN_US && width <= RC_VALID_MAX_US) {
            portENTER_CRITICAL_ISR(&mux_);
            pulseUs_[idx]     = static_cast<uint16_t>(width);
            lastValidUs_[idx] = now;
            portEXIT_CRITICAL_ISR(&mux_);
        }
    }
}
\end{lstlisting}

Les routines ISR sont placées en RAM (\texttt{IRAM\_ATTR}) pour éviter les
délais de lecture depuis le flash. La synchronisation avec la boucle principale
utilise \texttt{portMUX}, mécanisme FreeRTOS approprié pour les ESP32 dual-core.

\subsubsection{Toutes les consignes actionneurs en microsecondes}

Tous les actionneurs sont commandés exclusivement en microsecondes (µs),
jamais en degrés. Cela évite les conversions d'unités sources d'erreurs et
correspond directement au protocole PWM des servos et ESC.

\subsubsection{Slew rate — protection contre les à-coups}

\begin{lstlisting}[style=cpp, caption={McpwmActuators.cpp -- limitation de variation}]
uint16_t McpwmActuators::slew(uint16_t current, uint16_t target,
                               uint16_t step) {
    if (current < target) {
        uint32_t next = static_cast<uint32_t>(current) + step;
        return (next >= target) ? target : static_cast<uint16_t>(next);
    }
    if (current > target) {
        return (current > target + step) ?
               static_cast<uint16_t>(current - step) : target;
    }
    return current;
}
// Appele a chaque tick de controle (20 ms) avec ESC_SLEW_US = 30 uss/tick
outEsc1Us_ = slew(outEsc1Us_, e1t, Calibration::ESC_SLEW_US);
\end{lstlisting}

La variation maximale autorisée est de 30~µs par tick de contrôle (20~ms),
soit 1500~µs/s maximum. Cela protège la transmission mécanique contre les
changements brusques de consigne.

\subsubsection{Armement ESC — sécurité démarrage}

\begin{lstlisting}[style=cpp, caption={ManualController.cpp -- procedure d'armement}]
if (!escArmed_) {
    if (frame.ch3 <= Calibration::ESC_ARM_MAX_US) {
        if (armStartMs_ == 0) armStartMs_ = nowMs;
        if ((nowMs - armStartMs_) >= Calibration::ESC_ARM_MS) {
            escArmed_ = true;  // Arme apres 2 secondes en dessous du seuil
        }
    } else {
        armStartMs_ = 0;       // Reset si les gaz remontent
    }
    cmd.esc1Us = Calibration::ESC_STOP_US;
    return cmd;                // Pas de commande moteur avant armement
}
\end{lstlisting}

L'ESC ne peut être armé que si le canal de gaz reste sous le seuil
(\us{1300}) pendant 2~secondes consécutives. Tout dépassement remet le
compteur à zéro.

% =============================================================================
\section{Phase 1 — Contrôle manuel RC}
% =============================================================================

\subsection{Réalisation}

La Phase~1 porte sur le contrôle manuel complet par radiocommande.
Fichiers implémentés~:

\begin{longtable}{L{4.5cm} L{8.5cm}}
  \toprule
  \rowcolor{tableheadbg}
  \textcolor{tableheadfg}{\textbf{Fichier}} &
  \textcolor{tableheadfg}{\textbf{Rôle}} \\
  \midrule
  \endhead
  \bottomrule
  \endfoot
  \file{AxpPower.h/.cpp} & Init AXP192~: LDO2 (LoRa), LDO3 (GPS),
                           DCDC1 (3,3~V), ADC batterie activé \\[3pt]
  \file{RcReceiver.h/.cpp} & 4 canaux, ISR CHANGE, portMUX, timeout 100~ms \\[3pt]
  \file{McpwmActuators.h/.cpp} & MCPWM timer 0 (voile + safran),
                                  timer 1 (ESC), slew rate \\[3pt]
  \file{ModeManager.h/.cpp} & CH5 → Failsafe / ManualServo / ManualProp /
                               Automatic \\[3pt]
  \file{ManualController.h/.cpp} & Aileron binaire, safran proportionnel,
                                   armement + gaz ESC \\[3pt]
  \file{DroneApp.h/.cpp} & Orchestrateur~: ticks 20~ms (contrôle),
                           200~ms (debug), 1~s (LoRa) \\[3pt]
  \file{Types.h} & \texttt{ControlMode}, \texttt{RcFrame},
                   \texttt{ActuatorCommand} \\
\end{longtable}

\subsection{Modes de vol}

\begin{longtable}{C{2cm} C{3cm} L{8cm}}
  \toprule
  \rowcolor{tableheadbg}
  \textcolor{tableheadfg}{\textbf{CH5 (µs)}} &
  \textcolor{tableheadfg}{\textbf{Mode}} &
  \textcolor{tableheadfg}{\textbf{Comportement}} \\
  \midrule
  \endhead
  \bottomrule
  \endfoot
  $< 1000$ & Automatic & Navigation GPS autonome \\[3pt]
  1400--1600 & ManualServo & CH2 $\to$ aileron (bascule $\pm10^{\circ}$),
               CH4 $\to$ safran (proportionnel) \\[3pt]
  $> 1800$ & ManualProp & CH3 $\to$ gaz ESC (avec armement),
             CH4 $\to$ safran \\[3pt]
  Signal absent & Failsafe & Tous actionneurs aux positions de sécurité \\
\end{longtable}

\subsection{Contrôle de l'aileron — logique binaire}

L'aileron du drone n'a que deux positions utiles ($+10^{\circ}$ ou $-10^{\circ}$).
Le canal CH2 est lu de façon binaire~: si le manche dépasse la bande
morte ($\pm35$~µs autour de 1500~µs), l'état bascule et reste mémorisé
jusqu'au prochain changement~:

\begin{lstlisting}[style=cpp, caption={Logique aileron binaire avec etat memorise}]
if (!sailStateInitialized_) {
    sailState_ = (frame.ch2 >= RC_MID_US) ? +1 : -1;
    sailStateInitialized_ = true;
}
if (frame.ch2 > RC_MID_US + db) sailState_ = +1;
else if (frame.ch2 < RC_MID_US - db) sailState_ = -1;
// Pas de reset dans la bande morte : l'etat precedent est maintenu

cmd.sailUs = (sailState_ > 0) ? SAIL_PLUS_US : SAIL_MINUS_US;
\end{lstlisting}

\subsection{Résultats des tests — Phase 1}

\begin{itemize}
  \item \done\ 106 tests unitaires compilés et exécutés nativement sur Linux
    (sans matériel) — tous passent
  \item \done\ Compilation arduino-cli~: 27~\% flash, 7~\% RAM
  \item \done\ Aileron bascule correctement entre \us{1465} et \us{1575}
  \item \done\ Safran répond proportionnellement au manche CH4
  \item \done\ Armement ESC fonctionnel (2~s sous le seuil)
  \item \done\ Failsafe déclenché en moins de 100~ms après perte du signal
\end{itemize}

% =============================================================================
\section{Phase 2 — GPS et navigation}
% =============================================================================

\subsection{Réalisation}

\begin{longtable}{L{4.5cm} L{8.5cm}}
  \toprule
  \rowcolor{tableheadbg}
  \textcolor{tableheadfg}{\textbf{Fichier}} &
  \textcolor{tableheadfg}{\textbf{Rôle}} \\
  \midrule
  \endhead
  \bottomrule
  \endfoot
  \file{GpsUart.h/.cpp} & Serial1 GPIO34/12, TinyGPSPlus, capture NMEA,
                          CFG-CFG reset + CFG-ANT au boot \\[3pt]
  \file{Navigator.h/.cpp} & Haversine (distance en m + relèvement en degrés),
                            fonctions stateless dans un namespace \\[3pt]
  \file{MissionPlan.h} & Liste de 16 waypoints max + mode (Linéaire/Circuit) \\[3pt]
  \file{MissionManager.h/.cpp} & Machine à états~: Idle → Running →
                                  Returning → Complete \\
\end{longtable}

\subsection{Problème critique — GPS muet après config flash}

\begin{warningbox}
  \textbf{Symptôme~:} le compteur \texttt{chars} reste bloqué à 10 (un frame
  UBX ACK-ACK), le GPS est alimenté mais ne produit aucune phrase NMEA.

  \textbf{Cause~:} une session précédente avec l'IDE Arduino ou u-center avait
  sauvegardé en flash interne du NEO-6M une configuration CFG-PRT désactivant
  la sortie NMEA. Cette configuration survit aux coupures d'alimentation et au
  retrait de la batterie.

  \textbf{Solution~:} envoyer un reset usine UBX CFG-CFG au démarrage~:
\end{warningbox}

\begin{lstlisting}[style=cpp, caption={GpsUart.cpp -- reset usine NEO-6M au demarrage}]
// UBX CFG-CFG : efface la config flash, charge les valeurs ROM
static const uint8_t cfgReset[] = {
    0xB5,0x62, 0x06,0x09,         // en-tete UBX + classe/ID CFG-CFG
    0x0D,0x00,                    // longueur payload = 13 octets
    0xFF,0xFF,0x00,0x00,          // clearMask = 0xFFFF (effacer tout)
    0x00,0x00,0x00,0x00,          // saveMask  = 0x0000
    0xFF,0xFF,0x00,0x00,          // loadMask  = 0xFFFF (charger ROM)
    0x17,                         // deviceMask = SPI+I2C+Flash+EEPROM
    0x2F,0xAE                     // CRC
};
Serial1.write(cfgReset, sizeof(cfgReset));
delay(500);
\end{lstlisting}

\subsection{Problème critique — Antenne active non alimentée}

\begin{warningbox}
  \textbf{Symptôme~:} NMEA reçu (\texttt{chars} s'incrémente), mais
  \texttt{visible = 0} satellites même en extérieur sous ciel dégagé,
  avec antenne externe connectée.

  \textbf{Cause~:} le superviseur d'antenne du NEO-6M est désactivé par
  défaut. Sans l'activer~:
  \begin{enumerate}
    \item Aucune tension de polarisation n'est fournie au LNA de l'antenne
      active Taoglas~: le LNA est sans alimentation.
    \item Le commutateur RF du T-Beam (\texttt{ANT\_FLAG}) reste sur l'antenne
      céramique interne~: le connecteur u.FL est court-circuité.
  \end{enumerate}
\end{warningbox}

\begin{lstlisting}[style=cpp, caption={GpsUart.cpp -- activation superviseur d'antenne (CFG-ANT)}]
// UBX CFG-ANT : flags = 0x001B active le superviseur d'antenne
// (svcs=1: tension RF, scd=1: detection court-circuit,
//  pdwnOnSCD=1: coupure sur court-circuit, recovery=1)
static const uint8_t cfgAnt[] = {
    0xB5,0x62, 0x06,0x13,
    0x04,0x00,
    0x1B,0x00,          // flags = 0x001B
    0x00,0x00,          // pins = 0 (valeur par defaut)
    0x38,0x43
};
Serial1.write(cfgAnt, sizeof(cfgAnt));
delay(200);
\end{lstlisting}

\begin{okbox}
  Résultat après correction~: fix acquis en moins de 60~s en extérieur,
  7 satellites visibles, HDOP = 2,2, position confirmée (zone Brest,
  48.360476° N / $-$4.566822° W).
\end{okbox}

\subsection{Note sur le démarrage à froid (cold start)}

Sans batterie LiPo connectée, chaque mise sous tension perd l'almanach GPS
stocké. Le démarrage à froid prend environ 2 minutes en extérieur dégagé.
Avec la batterie, les éphémérides sont conservés et le fix s'obtient en
quelques secondes (démarrage chaud). La LED TIMEPULSE sur le NEO-6M clignote
à 1~Hz quand un fix est acquis.

% =============================================================================
\section{Phase 3 — LoRa et télémétrie}
% =============================================================================

\subsection{Réalisation}

\begin{longtable}{L{4.5cm} L{8.5cm}}
  \toprule
  \rowcolor{tableheadbg}
  \textcolor{tableheadfg}{\textbf{Fichier}} &
  \textcolor{tableheadfg}{\textbf{Rôle}} \\
  \midrule
  \endhead
  \bottomrule
  \endfoot
  \file{LoRaRadio.h/.cpp} & Driver SX1276, SPI HSPI, 433~MHz,
                            TX bloquant, RX par polling \\[3pt]
  \file{LoRaComm.h/.cpp} & Heartbeat JSON TX à 1~Hz, dispatcher de commandes
                           (navigate, stop, home, waypoints, wind, restart) \\[3pt]
  \file{transceiver.ino} & Sketch station sol~: CLI compacte + passerelle
                           JSON bidirectionnelle \\
\end{longtable}

\subsection{Protocole JSON LoRa}

\begin{lstlisting}[style=json, caption={Heartbeat drone vers station sol (p\'eriode 1 s)}]
{
  "origin": "boat",
  "type":   "info",
  "message": {
    "mode":         "navigate",
    "location":     [48.360476, -4.566822],
    "servos":       {"sail": 0, "rudder": 0},
    "control_mode": "autonomous",
    "heading":      270.5,
    "wind":         180,
    "bat":          7.42,
    "waypoints":    {"total": 3, "current": 1}
  }
}
\end{lstlisting}

\begin{lstlisting}[style=json, caption={Commandes station sol vers drone}]
{"origin":"server","type":"command","message":{"navigate":true}}
{"origin":"server","type":"command","message":{"stop":true}}
{"origin":"server","type":"command","message":{"restart":true}}
{"origin":"server","type":"command","message":
  {"waypoints":{"number":2,"points":"48.36,-4.56,48.37,-4.55"}}}
{"origin":"server","type":"command","message":
  {"home":{"lat":48.36,"lon":-4.56}}}
\end{lstlisting}

\subsection{Résultats}

\begin{itemize}
  \item \done\ TX confirmé (compteur \texttt{tx} s'incrémente à 1~Hz)
  \item \done\ 7 commandes testées via script Python (\texttt{lora\_uplink\_test.py})
  \item \done\ RSSI~: $-101$ à $-106$~dBm à 5~cm de distance (valeur attendue)
  \item \done\ GPS et LoRa opèrent simultanément sans interférence
  \item \todo\ Test de portée réelle en extérieur (prévu Phase~7)
\end{itemize}

\begin{infobox}[Note sur le TX bloquant]
  \texttt{LoRa.endPacket()} est synchrone. À SF7/BW125~kHz, un paquet
  de $\approx$210 octets nécessite $\approx$450~ms de transmission, bloquant
  la boucle principale une fois par seconde. Acceptable en Phase~3~; à passer
  en mode asynchrone avant la Phase~6 si le jitter du loop de contrôle
  devient critique.
\end{infobox}

% =============================================================================
\section{PCB Version 2 — Améliorations matérielles}
% =============================================================================

\subsection{Motivations}

La Version~2 du PCB corrige tous les problèmes identifiés sur la V1 et intègre
les retours d'expérience du développement logiciel~:

\begin{enumerate}
  \item Déplacement des broches RC vers des GPIO sans conflits
  \item Amplification de niveau logique 3,3~V → 5~V par transistors NPN BC548
  \item Diviseur résistif haute impédance pour la mesure de tension batterie
  \item Libération de GPIO21/22 pour une utilisation I2C future
  \item Suppression du CH6 (libère GPIO23 pour le RST LoRa)
  \item Un seul ESC (suppression ESC2)
\end{enumerate}

\subsection{Correspondance des broches — PCB V2}

\begin{longtable}{C{2.5cm} C{2cm} L{7cm}}
  \toprule
  \rowcolor{tableheadbg}
  \textcolor{tableheadfg}{\textbf{Signal}} &
  \textcolor{tableheadfg}{\textbf{GPIO}} &
  \textcolor{tableheadfg}{\textbf{Notes}} \\
  \midrule
  \endhead
  \bottomrule
  \endfoot
  RC CH2 (aileron) & 39 & SVN, entrée seule, pas de drain continu \\[3pt]
  RC CH3 (gaz) & 14 & Libéré (n'est plus partagé avec I2C) \\[3pt]
  RC CH4 (safran) & 13 & Libéré \\[3pt]
  RC CH5 (mode) & 4 & Commutateur 3 positions \\[3pt]
  Servo voile (PWM1) & 2 & Via transistor BC548 \\[3pt]
  Servo safran (PWM2) & 25 & Via transistor BC548 \\[3pt]
  ESC1 & 15 & Via transistor BC548 \\[3pt]
  ADC batterie & 36 & SVP, ADC1\_CH0, entrée seule \\[3pt]
  LoRa SCK & 5 & SPI HSPI \\[3pt]
  LoRa MISO & 19 & SPI HSPI \\[3pt]
  LoRa MOSI & 27 & SPI HSPI \\[3pt]
  LoRa CS & 18 & Actif bas \\[3pt]
  LoRa RST & 23 & Maintenant libre (CH6 supprimé) \\[3pt]
  LoRa IRQ (DIO0) & 26 & \\[3pt]
  GPS RX & 34 & Entrée seule \\[3pt]
  GPS TX & 12 & \\[3pt]
  I2C SDA & 21 & AXP192 interne + connecteur d'extension \\[3pt]
  I2C SCL & 22 & AXP192 interne + connecteur d'extension \\
\end{longtable}

\subsection{Amplification de niveau logique — Transistors BC548}

Les sorties PWM de l'ESP32 sont à 3,3~V. Les servos et ESC Pro-Tronik
attendent 5~V. La V2 utilise des transistors NPN BC548 en montage émetteur
commun avec une résistance de base 5,1~k$\Omega$ et une résistance de
collecteur 11,3~k$\Omega$~:

\begin{infobox}[Découverte importante — DUTY\_MODE\_0 pour les deux canaux]
  Initialement, le canal OPR\_B (safran) utilisait \texttt{MCPWM\_DUTY\_MODE\_1}
  (inversion), sous l'hypothèse que le transistor NPN inversait le signal et
  qu'il fallait pré-inverser. En réalité, \textbf{les deux canaux} (OPR\_A et
  OPR\_B) fonctionnent correctement avec \texttt{DUTY\_MODE\_0} à travers
  le circuit BC548. L'utilisation de \texttt{DUTY\_MODE\_1} sur OPR\_B
  produisait une double inversion, rendant le signal inutilisable pour le
  servo safran. La correction consiste à utiliser \texttt{DUTY\_MODE\_0}
  (actif haut) pour les deux canaux.
\end{infobox}

\subsection{Mesure de tension batterie — diviseur haute impédance}

\begin{lstlisting}[style=cpp, caption={BatteryAdc.cpp -- calcul de tension}]
// R5 = 562 kOhm (cote haut, batterie vers ADC)
// R6 = 120 kOhm (cote bas, vers GND)
// V_bat = V_adc x (R5 + R6) / R6 = V_adc x 5,683
static constexpr float DIVIDER_RATIO =
    (562000.0f + 120000.0f) / 120000.0f;  // = 5,683

void BatteryAdc::begin() {
    analogSetPinAttenuation(BoardConfig::BATTERY_ADC_PIN, ADC_11db);
}
float BatteryAdc::readVolts() {
    uint32_t vPin_mV = analogReadMilliVolts(BoardConfig::BATTERY_ADC_PIN);
    return (vPin_mV * DIVIDER_RATIO) / 1000.0f;
}
\end{lstlisting}

Avec des résistances de 562~k$\Omega$ et 120~k$\Omega$, la consommation du
diviseur à pleine charge (7,4~V) est de $7{,}4 / 682000 \approx 10{,}9$~µA —
négligeable.

\subsection{Problème rencontré — résistances inversées}

\begin{warningbox}
  \textbf{Symptôme~:} lecture ADC affichant 17,9~V au lieu de $\approx$5~V.

  \textbf{Cause~:} les résistances physiques étaient câblées à l'inverse~:
  120~k$\Omega$ côté haut et 562~k$\Omega$ côté bas. Le point de mesure ADC
  se retrouvait à $5{,}0 \times 562/(120+562) \approx 4{,}12$~V, saturant
  l'ADC. La valeur lue était multipliée par le ratio inversé
  $682/120 = 5{,}683$, d'où 17,9~V.

  \textbf{Solution~:} inversion physique des résistances. Le code était correct.
\end{warningbox}

\subsection{Gestion de l'AXP192 — I2C persistant}

Dans la V1 du code, \texttt{Wire.end()} était appelé à la fin de l'init
AXP192 pour libérer GPIO21/22 destinés au RC. Dans la V2, ces broches sont
libres~: GPIO21 et GPIO22 restent dédiés à l'I2C. \texttt{Wire.end()} est
supprimé du chemin de succès~; le bus I2C reste actif pour de futurs
capteurs via un connecteur d'extension.

% =============================================================================
\section{Problèmes rencontrés et solutions détaillées}
% =============================================================================

Cette section recense tous les problèmes significatifs rencontrés au cours
du développement, classés par domaine.

\subsection{Matériel — ESP32 / T-Beam}

\subsubsection{GPIO35 — drain de 20~mA continu}

\begin{warningbox}
  Le T-Beam V1.1 comporte un diviseur résistif d'impédance totale
  $\approx370$~$\Omega$ permanentement connecté à \gpio{35}. Ce diviseur
  consomme environ 20~mA en continu depuis la batterie, soit $\approx$1,8~W,
  indépendamment de toute lecture ADC.
\end{warningbox}

\begin{okbox}
  \textbf{Décision~:} ne jamais lire \gpio{35}. L'ADC batterie est implémenté
  sur \gpio{36} (SVP, ADC1, entrée seule, sans drain) avec diviseur externe
  haute impédance (562~k$\Omega$ + 120~k$\Omega$).
\end{okbox}

\subsubsection{LoRa RST — pin incorrect dans la documentation}

\begin{warningbox}
  La documentation de référence de LilyGO et divers exemples en ligne
  indiquent \gpio{14} pour le RESET du SX1276. Le hardware T-Beam V1.1
  câble physiquement le RESET sur \gpio{23}.
\end{warningbox}

\begin{okbox}
  Le code Post-Claude utilise \texttt{LORA\_RST\_PIN = 23} (correct).
  Dans la PCBv3/main actuelle, GPIO23 est libre (CH6 supprimé), donc le
  RESET matériel est maintenant fonctionnel.
\end{okbox}

\subsubsection{AXP192 — GPS non alimenté}

\begin{warningbox}
  Le GPS (u-blox NEO-6M) est alimenté par le rail LDO3 de l'AXP192.
  Sans activer ce rail en code, le GPS ne reçoit aucune alimentation et
  reste silencieux, même si le module est physiquement présent.
\end{warningbox}

\begin{okbox}
  \texttt{AxpPower::begin()} active explicitement LDO3 à 3300~mV avant
  toute initialisation du GPS. LDO2 (LoRa) et DCDC1 (3,3~V) sont
  également activés.
\end{okbox}

\subsection{Logiciel — Firmware}

\subsubsection{DUTY\_MODE\_1 sur OPR\_B — servo safran inopérant}

\begin{warningbox}
  \textbf{Symptôme~:} le servo safran (Regatta ECO II, GPIO25) ne répondait
  à aucune commande. Le servo aileron (GPIO2) fonctionnait correctement.
  Les valeurs de consigne semblaient correctes dans les logs.

  \textbf{Cause~:} le code utilisait \texttt{MCPWM\_DUTY\_MODE\_1} (signal
  inversé) sur OPR\_B (safran), supposant que le transistor NPN BC548
  inverserait le signal et que la pré-inversion était nécessaire.
  En réalité, le circuit BC548 en émetteur commun produit une inversion
  de son côté~; appliquer une inversion côté MCPWM entraîne une double
  inversion, résultant en un signal inutilisable (\us{19000} de haut,
  \us{1000} de bas au lieu de l'inverse).

  \textbf{Comportement observé~:} signal en sortie du transistor de
  fréquence 50~Hz mais avec un rapport cyclique inverse~: le servo interprète
  le temps bas comme la commande, ce qui sort de la plage valide.
\end{warningbox}

\begin{okbox}
  Suppression de l'appel \texttt{mcpwm\_set\_duty\_type()} pour OPR\_B~:
  les deux canaux utilisent maintenant \texttt{DUTY\_MODE\_0} (actif haut)
  sans inversion. Le servo safran répond immédiatement.
\end{okbox}

\subsubsection{Seuil d'armement ESC trop bas}

\begin{warningbox}
  \textbf{Symptôme~:} l'ESC ne s'arme jamais malgré la procédure correcte.
  Les logs indiquent~: \texttt{CH3=1265} alors que
  \texttt{ESC\_ARM\_MAX\_US=1050}.

  \textbf{Cause~:} l'émetteur Pro-Tronik PTR-6A n'atteint pas 1000~µs
  à la position minimale de son manche de gaz~: le minimum réel mesuré
  est $\approx1265$~µs. Le seuil d'armement de 1050~µs n'est donc jamais
  atteint physiquement.
\end{warningbox}

\begin{okbox}
  Le seuil \texttt{ESC\_ARM\_MAX\_US} est porté de 1050 à 1300~µs dans
  \file{Calibration.h}, permettant l'armement au minimum de manche réel
  tout en conservant une marge de sécurité suffisante.
\end{okbox}

\subsubsection{Mode bloqué en Automatic — confusion CH4/CH5}

\begin{warningbox}
  \textbf{Symptôme~:} après mise sous tension, le drone restait
  systématiquement en mode Automatic, même avec le commutateur de mode
  en position ManualServo.

  \textbf{Cause~:} dans la version initiale, \texttt{ModeManager} lisait
  CH4 pour déterminer le mode de vol, alors que le commutateur 3 positions
  physique est câblé sur CH5. Le canal CH4 (manche analogique) renvoyait
  une valeur $< 1000$~µs au repos, déclenchant le mode Automatic.
\end{warningbox}

\begin{okbox}
  Remplacement de toutes les références \texttt{frame.ch4} par
  \texttt{frame.ch5} dans \texttt{ModeManager.cpp}, et renommage des
  constantes \texttt{CH4\_*} en \texttt{CH5\_*} dans \texttt{BoardConfig.h}.
\end{okbox}

\subsubsection{GPS — config flash désactivant NMEA}

Voir Section~6.2 pour la description complète et la solution.

\subsubsection{GPS — antenne active non alimentée}

Voir Section~6.3 pour la description complète et la solution.

\subsubsection{Port série occupé lors du flash}

\begin{warningbox}
  \textbf{Symptôme~:} échec du transfert firmware avec l'erreur
  \texttt{could not open port /dev/ttyUSB0: [Errno 16] Device or resource busy}.

  \textbf{Cause~:} un processus (moniteur série ou script Python) tenait
  le port série ouvert pendant la tentative de flash.
\end{warningbox}

\begin{okbox}
  Identifier et tuer le processus~: \texttt{fuser /dev/ttyUSB0} puis
  \texttt{kill <PID>}, ou simplement fermer le moniteur série avant le flash.
\end{okbox}

% =============================================================================
\section{État actuel du système}
% =============================================================================

\subsection{Résumé des phases}

\begin{longtable}{C{1.5cm} L{5cm} C{2cm} L{4.5cm}}
  \toprule
  \rowcolor{tableheadbg}
  \textcolor{tableheadfg}{\textbf{Phase}} &
  \textcolor{tableheadfg}{\textbf{Objectif}} &
  \textcolor{tableheadfg}{\textbf{État}} &
  \textcolor{tableheadfg}{\textbf{Notes}} \\
  \midrule
  \endhead
  \bottomrule
  \endfoot
  1 & Contrôle manuel RC & \done\ Complet & 106 tests, testé en vol \\[3pt]
  2 & GPS + navigation & \done\ Complet & Fix testé en extérieur \\[3pt]
  3 & LoRa + télémétrie & \done\ Complet & TX/RX vérifiés \\[3pt]
  PCB V2 & Nouveau câblage + I2C libre & \done\ Complet & Compilé et testé \\[3pt]
  4 & Capteurs I2C & \todo\ À faire & GPIO21/22 disponibles \\[3pt]
  5 & Estimateur de vent & \todo\ À faire & 30~m + 2 amures \\[3pt]
  6 & Navigation autonome & \todo\ À faire & SailAutoMode \\[3pt]
  7 & Mission complète & \todo\ À faire & Waypoints + LoRa \\
\end{longtable}

\subsection{Utilisation flash / RAM}

\begin{longtable}{L{5cm} C{2.5cm} C{2.5cm}}
  \toprule
  \rowcolor{tableheadbg}
  \textcolor{tableheadfg}{\textbf{Sketch}} &
  \textcolor{tableheadfg}{\textbf{Flash}} &
  \textcolor{tableheadfg}{\textbf{RAM}} \\
  \midrule
  \endhead
  \bottomrule
  \endfoot
  \file{main/main.ino} (drone) & 27~\% & 7~\% \\[3pt]
  \file{transceiver/transceiver.ino} & 23~\% & 6~\% \\
\end{longtable}

\subsection{Valeurs de calibration actuelles}

\begin{longtable}{L{5cm} C{2cm} L{5.5cm}}
  \toprule
  \rowcolor{tableheadbg}
  \textcolor{tableheadfg}{\textbf{Paramètre}} &
  \textcolor{tableheadfg}{\textbf{Valeur}} &
  \textcolor{tableheadfg}{\textbf{Notes}} \\
  \midrule
  \endhead
  \bottomrule
  \endfoot
  \texttt{SAIL\_CENTER\_US} & 1520~µs & Centre Futaba S3003 (non 1500!) \\[3pt]
  \texttt{SAIL\_PLUS\_US} & 1575~µs & $+10^{\circ}$ — à vérifier sur banc \\[3pt]
  \texttt{SAIL\_MINUS\_US} & 1465~µs & $-10^{\circ}$ — à vérifier sur banc \\[3pt]
  \texttt{ROTOR\_CENTER\_US} & 1500~µs & Safran centré (position de repos) \\[3pt]
  \texttt{ROTOR\_MIN\_US} & 1000~µs & Butée basse \\[3pt]
  \texttt{ROTOR\_MAX\_US} & 2000~µs & Butée haute \\[3pt]
  \texttt{ESC\_STOP\_US} & 1000~µs & Moteur arrêté \\[3pt]
  \texttt{ESC\_MAX\_US} & 2000~µs & Plein régime \\[3pt]
  \texttt{ESC\_ARM\_MAX\_US} & 1300~µs & Seuil d'armement \\[3pt]
  \texttt{ESC\_ARM\_MS} & 2000~ms & Durée maintien pour armer \\[3pt]
  \texttt{ESC\_SLEW\_US} & 30~µs/tick & Limitation variation ESC \\[3pt]
  \texttt{RC\_DEADBAND\_US} & 35~µs & Bande morte manche \\[3pt]
  Ratio diviseur batterie & 5,683 & (562k + 120k) / 120k \\
\end{longtable}

% =============================================================================
\section{Plan de développement futur}
% =============================================================================

\subsection{Phase 4 — Capteurs I2C}

Maintenant que GPIO21 et GPIO22 sont libres (bus I2C persistant, plus de
\texttt{Wire.end()}), des capteurs peuvent être connectés au connecteur
d'extension du PCB V2. Travaux à effectuer~:

\begin{itemize}
  \item Identifier les capteurs requis et leurs adresses I2C
  \item Implémenter \file{SensorBus.cpp} avec scan de bus et collecte d'échantillons
  \item Implémenter \file{SensorManager} (polling à 500~ms, buffering)
\end{itemize}

\subsection{Phase 5 — Estimateur de vent}

Sans capteur de vent, la direction est inférée depuis le track GPS~:

\begin{enumerate}
  \item Collecter le cap GPS sur deux amures différentes (état aileron $\pm1$)
  \item Le vent bissecte l'angle entre les deux caps
  \item Minimum 30~m parcourus sur chaque amure pour valider l'estimation
  \item Possibilité de forcer la direction via commande LoRa (\texttt{wind-command})
\end{enumerate}

\subsection{Phase 6 — Navigation autonome à la voile}

\begin{lstlisting}[style=cpp, caption={Algorithme SailAutoMode (pseudo-code)}]
// Entrees : position GPS, waypoint, cap GPS, direction vent estimee
// Sorties : sailUs (+/- SAIL_DELTA), rotorUs (position continue)
float bearing = Navigator::bearingDeg(pos, waypoint);
float relWind = normalizeAngle(bearing - windDir);
if (abs(relWind) < UPWIND_ZONE_DEG) {
    // Waypoint dans la zone morte (vent debout) -- virement de bord
    // Calculer la route de louvoyage optimale
} else {
    // Route directe possible -- choisir l'amure et regler le cap
    sailState = chooseTack(relWind);
    cmd.sailUs = sailState > 0 ? SAIL_PLUS_US : SAIL_MINUS_US;
    cmd.rotorUs = steerToward(pos, waypoint, heading);
}
\end{lstlisting}

\begin{warningbox}
  Un \textbf{WinchTracker} logiciel sera nécessaire~: le Regatta ECO II
  n'a pas de retour de position. Il faudra intégrer vitesse × temps pour
  estimer la position du treuil et détecter les butées mécaniques.
\end{warningbox}

\subsection{Phase 7 — Mission complète}

\begin{itemize}
  \item Chargement des waypoints par LoRa depuis la station sol
  \item Démarrage de la mission, monitoring par heartbeat LoRa
  \item Journalisation des données capteurs sur SPIFFS
  \item Test de portée LoRa en conditions réelles
\end{itemize}

% =============================================================================
\section{Environnement de développement}
% =============================================================================

\subsection{Outils}

\begin{longtable}{L{4cm} L{9cm}}
  \toprule
  \rowcolor{tableheadbg}
  \textcolor{tableheadfg}{\textbf{Outil}} &
  \textcolor{tableheadfg}{\textbf{Notes}} \\
  \midrule
  \endhead
  \bottomrule
  \endfoot
  Arduino IDE 2.3.8 & AppImage Linux, écriture code + upload + moniteur série \\[3pt]
  arduino-cli v1.4.1 & Compilation en ligne de commande (CI) \\[3pt]
  KiCad 9.0 & Conception PCB V1 et V2 \\[3pt]
  \texttt{pdflatex} & Compilation de ce rapport \\[3pt]
  ESP32 Arduino Core 3.x & Bibliothèque cible \\[3pt]
  AXP202X\_Library 1.1.2 & Gestion AXP192 \\[3pt]
  TinyGPSPlus 1.0.3 & Parseur NMEA \\[3pt]
  LoRa 0.8.0 & Driver SX1276 \\
\end{longtable}

\subsection{Commandes utiles}

\begin{lstlisting}[style=shell, caption={Compiler et flasher le firmware}]
# Compilation
~/bin/arduino-cli compile --fqbn esp32:esp32:t-beam \
    /chemin/vers/Post-Claude/main

# Flash (remplacer /dev/ttyUSB0 selon le systeme)
~/bin/arduino-cli upload --fqbn esp32:esp32:t-beam \
    --port /dev/ttyUSB0 /chemin/vers/Post-Claude/main

# Surveiller le port serie
python3 -m serial.tools.miniterm /dev/ttyUSB0 115200
\end{lstlisting}

% =============================================================================
\section{Journal de développement}
% =============================================================================

\begin{longtable}{C{2.5cm} L{10cm}}
  \toprule
  \rowcolor{tableheadbg}
  \textcolor{tableheadfg}{\textbf{Date}} &
  \textcolor{tableheadfg}{\textbf{Activité}} \\
  \midrule
  \endhead
  \bottomrule
  \endfoot
  2026-04-27 & Lecture du projet, analyse de l'architecture, création de CLAUDE.md.
              Recherche composants (Futaba S3003, Regatta ECO II, Pro-Tronik,
              T-Beam V1.1). Phase 1 codée (10 fichiers). Diagnostic drain GPIO35
              ($\approx$20~mA). \\[3pt]
  2026-04-27 & Tests unitaires natifs (59 tests). Bug ESC : \texttt{disarm()} manquant
              lors du passage en mode ManualServo — corrigé. Tous les tests passent. \\[3pt]
  2026-04-28 & 47 nouveaux tests (106 au total). Correction structure Arduino IDE
              (\texttt{src/} requis). Compilation arduino-cli : 24~\% flash, 7~\% RAM. \\[3pt]
  2026-04-28 & Phase 2 — GPS~: GpsUart, Navigator (haversine), MissionPlan,
              MissionManager. Diagnostic GPS froid en extérieur (sans LiPo = cold start). \\[3pt]
  2026-05-04 & Phase 3 — LoRa~: driver, LoRaComm (JSON), transceiver sol.
              TX vérifié. 7 commandes testées via Python. Phase 3 complète. \\[3pt]
  2026-05-20 & Correction MCPWM DUTY\_MODE\_1 → DUTY\_MODE\_0 sur OPR\_B (safran).
              Diagnostic et correction GPS (CFG-CFG + CFG-ANT). Fix GPS obtenu en extérieur. \\[3pt]
  2026-05-21 & Analyse schéma PCB V2 (KiCad). Validation des nouvelles connexions
              (BC548, diviseur, GPIO libres). \\[3pt]
  2026-05-22 & Refactoring PCBv3 (sketch test) : suppression ESC2, CH5 pour mode,
              CH4 pour safran. Correction seuil armement (1050 → 1300 µs).
              Compilation, flash, tests physiques. Aileron et safran fonctionnels.
              Armement ESC validé. \\[3pt]
  2026-05-24 & Mise à jour du sketch production (\texttt{main/}) pour le nouveau
              câblage PCB V2~: GPIO39/14/13/4 pour RC, GPIO15 pour ESC, GPIO36
              pour batterie, LoRa RST=23, Wire.end() supprimé.
              Compilation propre (27~\% flash, 7~\% RAM). \\
\end{longtable}

% =============================================================================
\appendix
\section{Spécifications des composants}
% =============================================================================

\subsection{Futaba S3003 — Servo aileron}

\begin{longtable}{L{5cm} L{7cm}}
  \toprule
  \rowcolor{tableheadbg}
  \textcolor{tableheadfg}{\textbf{Paramètre}} &
  \textcolor{tableheadfg}{\textbf{Valeur}} \\
  \midrule
  \endhead
  \bottomrule
  \endfoot
  Fréquence PWM & 50~Hz \\[3pt]
  Plage d'impulsions & 1000--2000~µs \\[3pt]
  Centre (neutre) & \textbf{1520~µs} (standard Futaba, pas 1500~!) \\[3pt]
  Bande morte & $\approx$100~µs \\[3pt]
  Tension d'alimentation & 4,8~V ou 6,0~V \\[3pt]
  Couple (4,8~V / 6~V) & 3,2~kg·cm / 4,1~kg·cm \\[3pt]
  Vitesse (4,8~V / 6~V) & 0,23~s/60° / 0,19~s/60° \\[3pt]
  Dimensions & 40,4 × 19,8 × 36~mm \\[3pt]
  Masse & 37,2~g \\
\end{longtable}

\subsection{Graupner Regatta ECO II (mod. 5176) — Treuil safran}

\begin{longtable}{L{5cm} L{7cm}}
  \toprule
  \rowcolor{tableheadbg}
  \textcolor{tableheadfg}{\textbf{Paramètre}} &
  \textcolor{tableheadfg}{\textbf{Valeur}} \\
  \midrule
  \endhead
  \bottomrule
  \endfoot
  Type & Treuil multi-tour \textbf{positionnel} \\[3pt]
  Fréquence PWM & 50~Hz \\[3pt]
  Plage d'impulsions & 1000--2000~µs \\[3pt]
  Centre (repos) & 1500~µs (tient la position) \\[3pt]
  Course mécanique & 6 tours ($\pm$3 tours depuis le centre) \\[3pt]
  Tension & 4,8--7,2~V \\[3pt]
  Courant repos / charge & 300--350~mA / 2500--3300~mA \\[3pt]
  Couple (6~V / 7,4~V) & $\approx$9,5~kg·cm / $\approx$10,9~kg·cm \\[3pt]
  Dimensions & 40,6 × 20 × 38,9~mm \\[3pt]
  Masse & 56~g \\
\end{longtable}

\subsection{Pro-Tronik PTR-6A / R8X}

\begin{longtable}{L{5cm} L{7cm}}
  \toprule
  \rowcolor{tableheadbg}
  \textcolor{tableheadfg}{\textbf{Paramètre}} &
  \textcolor{tableheadfg}{\textbf{Valeur}} \\
  \midrule
  \endhead
  \bottomrule
  \endfoot
  Protocole & FHSS 2,4~GHz (propriétaire Pro-Tronik) \\[3pt]
  Canaux TX & 6 (4 proportionnels + CH5 3 positions + CH6 2 positions) \\[3pt]
  Canaux RX & 8 (R8X) \\[3pt]
  PWM neutre & 1500~µs \\[3pt]
  Plage PWM & 1000--2000~µs \\[3pt]
  Fréquence PWM & 50~Hz \\[3pt]
  Portée & $\approx$2400~m \\[3pt]
  Minimum manche gaz (mesuré) & $\approx$1265~µs (pas 1000~µs !) \\
\end{longtable}

\end{document}
